\documentclass{../templateReading/cdreading}

\newcommand{\authorinfo}{THE FORUM CENTER - NGUYỄN HOÀNG HUY}
\newcommand{\documenttitle}{TÀI LIỆU CHUYÊN ĐỀ READING}
\newcommand{\collectiontitle}{Level 3 — W5 (Jan 26-Feb 1)}

\begin{document}

\thispagestyle{firstpage}
\makeworkshopheader{\authorinfo}{\documenttitle}{\collectiontitle}
\pagestyle{otherpages}

\cdsection{Wealth in a cold climate}

\textbf{A}\par

Dr William Masters was reading a book about mosquitoes when inspiration struck. “There was this anecdote about the great yellow fever epidemic that hit Philadelphia in 1793,” Masters recalls. “This epidemic decimated the city until the first frost came.” The inclement weather froze out the insects, allowing Philadelphia to recover.

\textbf{B}\par

If weather could be the key to a city’s fortunes, Masters thought, then why not to the historical fortunes of nations? And could frost lie at the heart of one of the most enduring economic mysteries of all — why are almost all the wealthy, industrialised nations to be found at latitudes above 40 degrees? After two years of research, he thinks that he has found a piece of the puzzle. Masters, an agricultural economist from Purdue University in Indiana, and Margaret McMillan at Tufts University, Boston, show that annual frosts are among the factors that distinguish rich nations from poor ones. Their study is published this month in the Journal of Economic Growth. The pair speculates that cold snaps have two main benefits — they freeze pests that would otherwise destroy crops, and also freeze organisms, such as mosquitoes, that carry disease. The result is agricultural abundance a big workforce.

\textbf{C}\par

The academics took two sets of information. The first was average income for countries, the second climate data from the University of East Anglia. They found a curious tally between the sets. Countries having five or more frosty days a month are uniformly rich; those with fewer than five are impoverished. The authors speculate that the five-day figure is important; it could be the minimum time needed to kill pests in the soil. Masters says: “For example, Finland is a small country that is growing quickly, but Bolivia is a small country that isn’t growing at all. Perhaps climate has something to do with that.” In fact, limited frosts bring huge benefits to farmers. The chills kill insects or render them inactive; cold weather slows the break-up of plant and animal material in the soil, allowing it to become richer; and frosts ensure a build-up of moisture in the ground for spring, reducing dependence on seasonal rains. There are exceptions to the “cold equals rich” argument. There are well-heeled tropical countries such as Hong Kong and Singapore (both city-states, Masters notes), a result of their superior trading positions. Likewise, not all European countries axe moneyed — in the former communist colonies, economic potential was crushed by politics.

\textbf{D}\par

Masters stresses that climate will never be the overriding factor 一 the wealth of nations is too complicated to be attributable to just one factor. Climate, he feels, somehow combines with other factors — such as the presence of institutions, including governments, and access to trading routes — to determine whether a country will do well. Traditionally, Masters says, economists thought that institutions had the biggest effect on the economy, because they brought order to a country in the form of, for example, laws and property rights. With order, so the thinking went, came affluence. “But there are some problems that even countries with institutions have not been able to get around,” he says. “My feeling is that, as countries get richer, they get better institutions. And the accumulation of wealth and improvement in governing institutions are both helped by a favourable environment, including climate.”

\textbf{E}\par

This does not mean, he insists, that tropical countries are beyond economic help and destined to remain penniless. Instead, richer countries should change the way in which foreign aid is given. Instead of aid being geared towards improving governance, it should be spent on technology to improve agriculture and to combat disease. Masters cites one example: “There are regions in India that have been provided with irrigation — agricultural productivity has gone up and there has been an improvement in health.” Supplying vaccines against tropical diseases and developing crop varieties that can grow in the tropics would break the poverty cycle.

\textbf{F}\par

Other minds have applied themselves to the split between poor and rich nations, citing anthropological, climatic and zoological reasons for why temperate nations are the most affluent. In 350BC, Aristotle observed that “those who live in a cold climate… are full of spirit”. Jared Diamond, from the University of California at Los Angeles, pointed out in his book Guns, Germs and Steel that Eurasia is broadly aligned east-west, while Africa and the Americas are aligned north-south. So, in Europe, crops can spread quickly across latitudes because climates are similar. One of the first domesticated crops, einkorn wheat, spread quickly from the Middle East into Europe; it took twice as long for corn to spread from Mexico to what is now the eastern United States. This easy movement along similar latitudes in Eurasia would also have meant a faster dissemination of other technologies such as the wheel and writing, Diamond speculates. The region also boasted domesticated livestock, which could provide meat, wool and motive power in the fields. Blessed with such natural advantages, Eurasia was bound to take off economically.

\textbf{G}\par

John Gallup and Jeffrey Sachs, two US economists, have also pointed out striking correlations between the geographical location of countries and their wealth. They note that tropical countries between 23.45 degrees north and south of the equator are nearly all poor. In an article for the Harvard International Review, they concluded that “development surely seems to favour the temperate-zone economies, especially those in the northern hemisphere, and those that have managed to avoid both socialism and the ravages of war”. But Masters cautions against geographical determinism, the idea that tropical countries are beyond hope: “Human health and agriculture can be made better through scientific and technological research,” he says, “so we shouldn’t be writing off these countries. Take Singapore: without air conditioning, it wouldn’t be rich.”

\cdsection{Questions 14-20}

The reading passage has seven paragraphs, A-G.

\cdnoteinline{Choose the correct heading for paragraphs A-G from the list below.}

\cdnoteinline{Write the correct number, i-x, in boxes 14-20 on your answer sheet.}

List of Headings

\cdblue{\textbf{i}} The positive correlation between climate and wealth

\cdblue{\textbf{ii}} Other factors besides climate that influence wealth

\cdblue{\textbf{iii}} Inspiration from reading a book

\cdblue{\textbf{iv}} Other researchers’ results do not rule out exceptional cases

\cdblue{\textbf{v}} different attributes between Eurasia and Africa

\cdblue{\textbf{vi}} Low temperature benefits people and crops

\cdblue{\textbf{vii}} The importance of institution in traditional views.

\cdblue{\textbf{viii}} The spread of crops in Europe, Asia and other places

\cdblue{\textbf{ix}} The best way to use aid

\cdblue{\textbf{x}} confusions and exceptional

\textbf{14} Paragraph A

\textbf{15} Paragraph B

\textbf{16} Paragraph C

\textbf{17} Paragraph D

\textbf{18} Paragraph E

\textbf{19} Paragraph F

\textbf{20} Paragraph G

\cdsection{Questions 21-26}

\cdnoteinline{Complete the following summary of the paragraphs of Reading Passage}

Using NO MORE THAN TWO WORDS from the Reading Passage for each answer.

\cdnoteinline{Write your answers in boxes 21-26 on your answer sheet.}

Dr William Master read a book saying that a (an) 21………………….. which struck an American city of years ago was terminated by a cold frost. And academics found that there is a connection between climate and country’s wealthy as in the rich but small country of 22…………………..; Yet besides excellent surroundings and climate, one country still need to improve both their 23………………….. to achieve long prosperity,

Thanks to resembling weather condition across latitude in the continent of 24………………….. ’crops such as 25…………………… is bound to spread faster than from South America to the North. Other researchers also noted that even though geographical factors are important, a tropical country such as 26………………….. still became rich due to scientific advancement.

\cdblue{\textbf{Hướng dẫn giải}}

\cdsection{Phần I: Câu hỏi 14-20 (Matching Headings)}

Mục tiêu: Nắm bắt ý chính của đoạn văn và tìm tiêu đề tương ứng trong danh sách cho sẵn.

\cdstep{1. Chiến lược tiếp cận}
\begin{itemize}
\item Bước 1: Giải mã Danh sách Tiêu đề (List of Headings)
\begin{itemize}
\item (i) Mối tương quan tích cực giữa khí hậu và sự giàu có.
\item (ii) Các yếu tố khác ngoài khí hậu ảnh hưởng đến sự giàu có.
\item (iii) Cảm hứng từ việc đọc một cuốn sách.
\item (iv) Kết quả của các nhà nghiên cứu khác không loại trừ các trường hợp ngoại lệ.
\item (v) Các đặc điểm khác nhau giữa Á-Âu (Eurasia) và Châu Phi.
\item (vi) Nhiệt độ thấp mang lại lợi ích cho con người và mùa màng.
\item (vii) Tầm quan trọng của thể chế theo quan điểm truyền thống.
\item (viii) Sự lan rộng của cây trồng ở Châu Âu, Châu Á và những nơi khác.
\item (ix) Cách tốt nhất để sử dụng viện trợ.
\item (x) Những sự nhầm lẫn và ngoại lệ.
\end{itemize}
\item Bước 2: Đọc lướt và Đối chiếu
\begin{itemize}
\item Tìm từ khóa trong đoạn văn và so sánh với ý nghĩa của các tiêu đề.
\end{itemize}
\end{itemize}

\cdstep{2. Phân tích chi tiết từng đoạn văn}

Câu 14: Paragraph A
\begin{itemize}
\item Nội dung chính: Đoạn văn kể về việc Dr William Masters đang đọc sách thì nảy ra ý tưởng.
\item Từ khóa tìm trong bài: reading a book, inspiration struck, anecdote.
\item Từ khóa mấu chốt: "Dr William Masters was reading a book... when inspiration struck."
\item Phân tích \& Vị trí: Tác giả mô tả khoảnh khắc Masters tìm thấy cảm hứng ("inspiration struck") khi đang đọc sách ("reading a book"). Điều này khớp hoàn toàn với tiêu đề (iii).
\item Đáp án: iii
\end{itemize}

Câu 15: Paragraph B
\begin{itemize}
\item Nội dung chính: Đoạn văn giải thích lợi ích của thời tiết lạnh đối với nông nghiệp và sức khỏe (diệt sâu bệnh).
\item Từ khóa tìm trong bài: freeze pests, destroy crops, freeze organisms, carry disease, agricultural abundance.
\item Từ khóa mấu chốt: "...cold snaps have two main benefits — they freeze pests... and also freeze organisms... The result is agricultural abundance..."
\item Phân tích \& Vị trí: Đoạn văn liệt kê hai lợi ích chính ("two main benefits") của các đợt rét ("cold snaps"): đóng băng sâu bệnh hại mùa màng và sinh vật gây bệnh. Kết quả là sự phong phú về nông nghiệp ("agricultural abundance"). Điều này khớp với tiêu đề (vi) về lợi ích của nhiệt độ thấp.
\item Đáp án: vi
\end{itemize}

Câu 16: Paragraph C
\begin{itemize}
\item Nội dung chính: Đoạn văn đưa ra dữ liệu so sánh giữa số ngày sương giá và mức độ giàu có của quốc gia.
\item Từ khóa tìm trong bài: tally between the sets, five or more frosty days... uniformly rich, fewer than five are impoverished.
\item Từ khóa mấu chốt: "Countries having five or more frosty days a month are uniformly rich; those with fewer than five are impoverished."
\item Phân tích \& Vị trí: Tác giả tìm thấy mối liên hệ ("tally") rằng các nước có nhiều ngày sương giá thì giàu ("rich"), còn ít ngày sương giá thì nghèo ("impoverished"). Đây chính là mối tương quan tích cực giữa khí hậu và sự giàu có (tiêu đề i).
\item Đáp án: i
\end{itemize}

Câu 17: Paragraph D
\begin{itemize}
\item Nội dung chính: Masters nhấn mạnh rằng khí hậu không phải là yếu tố duy nhất, mà còn cần các yếu tố khác như thể chế.
\item Từ khóa tìm trong bài: climate will never be the overriding factor, combines with other factors, institutions.
\item Từ khóa mấu chốt: "Masters stresses that climate will never be the overriding factor... Climate, he feels, somehow combines with other factors — such as the presence of institutions..."
\item Phân tích \& Vị trí: Đoạn văn khẳng định khí hậu kết hợp với "other factors" (các yếu tố khác) như "institutions" (thể chế) để quyết định sự thịnh vượng. Điều này khớp với tiêu đề (ii).
\item Đáp án: ii
\end{itemize}

Câu 18: Paragraph E
\begin{itemize}
\item Nội dung chính: Đề xuất thay đổi cách viện trợ cho các nước nghèo (tập trung vào công nghệ thay vì quản trị).
\item Từ khóa tìm trong bài: richer countries should change the way in which foreign aid is given, spent on technology.
\item Từ khóa mấu chốt: "Instead of aid being geared towards improving governance, it should be spent on technology to improve agriculture..."
\item Phân tích \& Vị trí: Tác giả bàn về cách các nước giàu nên thay đổi phương thức viện trợ nước ngoài ("foreign aid"), cụ thể là nên chi cho công nghệ. Đây là đề xuất về cách sử dụng viện trợ tốt nhất (tiêu đề ix).
\item Đáp án: ix
\end{itemize}

Câu 19: Paragraph F
\begin{itemize}
\item Nội dung chính: So sánh sự khác biệt về trục địa lý giữa Á-Âu (Đông-Tây) và Châu Phi/Châu Mỹ (Bắc-Nam) ảnh hưởng đến sự lan truyền cây trồng.
\item Từ khóa tìm trong bài: Eurasia is broadly aligned east-west, Africa and the Americas are aligned north-south.
\item Từ khóa mấu chốt: "...Eurasia is broadly aligned east-west, while Africa and the Americas are aligned north-south."
\item Phân tích \& Vị trí: Đoạn văn chỉ ra sự khác biệt về thuộc tính địa lý ("aligned east-west" vs "aligned north-south") giữa Eurasia và Africa/Americas. Điều này khớp với tiêu đề (v).
\item Đáp án: v
\end{itemize}

Câu 20: Paragraph G
\begin{itemize}
\item Nội dung chính: Nhắc đến nghiên cứu của các nhà kinh tế khác về các nước nhiệt đới nghèo, nhưng Masters chỉ ra ngoại lệ (Singapore).
\item Từ khóa tìm trong bài: tropical countries... nearly all poor, Singapore... wouldn't be rich.
\item Từ khóa mấu chốt: "John Gallup and Jeffrey Sachs... pointed out striking correlations... But Masters cautions... Take Singapore..."
\item Phân tích \& Vị trí: Đoạn văn nhắc đến kết quả của "John Gallup and Jeffrey Sachs" (các nhà nghiên cứu khác) về việc các nước nhiệt đới thường nghèo. Tuy nhiên, Masters đưa ra ví dụ về Singapore như một ngoại lệ nhờ công nghệ. Điều này khớp với tiêu đề (iv) - kết quả của người khác không loại trừ ngoại lệ.
\item Đáp án: iv
\end{itemize}

\cdsection{Phần II: Câu hỏi 21-26 (Summary Completion)}

\cdnoteinline{Yêu cầu: Hoàn thành bản tóm tắt sử dụng KHÔNG QUÁ HAI TỪ từ bài đọc.}

Câu 21

Nội dung: Dr William Master read a book saying that a (an) 21………………….. which struck an American city of years ago was terminated by a cold frost.
\begin{enumerate}
\item Từ khóa tìm trong bài
\begin{itemize}
\item read a book
\item struck
\item American city
\item terminated by a cold frost
\end{itemize}
\item Từ khóa mấu chốt
\begin{itemize}
\item Đoạn A: "...anecdote about the great yellow fever epidemic that hit Philadelphia... The inclement weather froze out the insects..."
\end{itemize}
\item Phân tích \& Vị trí
\begin{itemize}
\item "Hit" tương đương với "struck". "Philadelphia" là một "American city". Sự kiện được nhắc đến là "yellow fever epidemic".
\end{itemize}
\item Đáp án: yellow-fever epidemic (hoặc epidemic)
\end{enumerate}

Câu 22

Nội dung: And academics found that there is a connection between climate and country’s wealthy as in the rich but small country of 22…………………..
\begin{enumerate}
\item Từ khóa tìm trong bài
\begin{itemize}
\item connection between climate and country’s wealthy
\item rich but small country
\end{itemize}
\item Từ khóa mấu chốt
\begin{itemize}
\item Đoạn C: "For example, Finland is a small country that is growing quickly..."
\end{itemize}
\item Phân tích \& Vị trí
\begin{itemize}
\item Đoạn văn lấy ví dụ về một quốc gia nhỏ ("small country") đang phát triển nhanh (giàu có) để minh họa cho mối liên hệ khí hậu. Đó là Phần Lan.
\end{itemize}
\item Đáp án: Finland
\end{enumerate}

Câu 23

Nội dung: Yet besides excellent surroundings and climate, one country still need to improve both their 23………………….. to achieve long prosperity,
\begin{enumerate}
\item Từ khóa tìm trong bài
\begin{itemize}
\item besides excellent surroundings and climate
\item improve both their ...
\item achieve long prosperity (wealth)
\end{itemize}
\item Từ khóa mấu chốt
\begin{itemize}
\item Đoạn D: "...accumulation of wealth and improvement in governing institutions are both helped by a favourable environment..."
\end{itemize}
\item Phân tích \& Vị trí
\begin{itemize}
\item Bài đọc nói rằng sự tích lũy của cải ("wealth" \~{} prosperity) và sự cải thiện trong "governing institutions" đều được hỗ trợ bởi môi trường. Do đó, yếu tố cần cải thiện là "governing institutions".
\end{itemize}
\item Đáp án: Governing institutions
\end{enumerate}

Câu 24

Nội dung: Thanks to resembling weather condition across latitude in the continent of 24…………………..
\begin{enumerate}
\item Từ khóa tìm trong bài
\begin{itemize}
\item resembling weather condition (climates are similar)
\item across latitude
\item continent of ...
\end{itemize}
\item Từ khóa mấu chốt
\begin{itemize}
\item Đoạn F: "So, in Europe, crops can spread quickly across latitudes because climates are similar."
\end{itemize}
\item Phân tích \& Vị trí
\begin{itemize}
\item "Climates are similar" tương ứng với "resembling weather condition". Đoạn văn nói điều này xảy ra ở châu lục/khu vực là Châu Âu ("Europe").
\end{itemize}
\item Đáp án: Europe
\end{enumerate}

Câu 25

Nội dung: ’crops such as 25…………………… is bound to spread faster than from South America to the North.
\begin{enumerate}
\item Từ khóa tìm trong bài
\begin{itemize}
\item crops such as ...
\item spread faster (quickly)
\end{itemize}
\item Từ khóa mấu chốt
\begin{itemize}
\item Đoạn F: "One of the first domesticated crops, einkorn wheat, spread quickly from the Middle East into Europe..."
\end{itemize}
\item Phân tích \& Vị trí
\begin{itemize}
\item Loại cây trồng ("crop") được nhắc đến là lan truyền nhanh ("spread quickly") chính là lúa mì einkorn.
\end{itemize}
\item Đáp án: einkorn Wheat
\end{enumerate}

Câu 26

Nội dung: Other researchers also noted that even though geographical factors are important, a tropical country such as 26………………….. still became rich due to scientific advancement.
\begin{enumerate}
\item Từ khóa tìm trong bài
\begin{itemize}
\item tropical country
\item became rich
\item scientific advancement (air conditioning)
\end{itemize}
\item Từ khóa mấu chốt
\begin{itemize}
\item Đoạn G: "Take Singapore: without air conditioning, it wouldn’t be rich."
\end{itemize}
\item Phân tích \& Vị trí
\begin{itemize}
\item Singapore là ví dụ về một quốc gia nhiệt đới ("tropical country") trở nên giàu có nhờ vào "air conditioning" (sản phẩm của tiến bộ khoa học kỹ thuật).
\end{itemize}
\item Đáp án: Singapore
\end{enumerate}

\end{document}
